\documentclass{situationist-preprint}

\title{The Standard of Taste}
\msid{AES-03-a}
\msfamily{AES-03}
\msdiscipline{Aesthetics}
\mstype{Principal manuscript}
\mskeywords{standard of taste; Hume; ideal critics; true judges; convergence; aesthetic disagreement; relativism; test of time}

\begin{document}
\maketitle

\begin{abstract}
Are there facts about taste? Hume held that sentiment is never mistaken, and yet that some judgments of taste are absurd; he reconciled the two by locating a standard of taste in the joint verdict of true judges, critics of strong sense, delicate sentiment, practice, comparison and freedom from prejudice. The standard is satisfiable only if such judges can be specified without circularity, can be approximated by actual critics, and agree. The last condition has been widely doubted, by Hume's own concession that temperament and the manners of an age produce blameless variation, and by later arguments that equally ideal sensibilities diverge. We argue that the verdicts of ideal critics converge. The key move is to separate verdicts of merit from preferences: every source of divergence among judges is either removed by one of Hume's qualifications, which we show must include a moral sensitivity contained in good sense, or confined to preferences among works whose merit differs by less than the influence of temperament, on which the joint verdict is one of parity. We prove a bound on the disagreement of imperfect critics that explains why actual critics approximate the standard, and we defend the result against idealization, sociological and relativist objections.
\end{abstract}

\section{The problem}

Two commonplaces about taste pull against each other. The first is that there is no disputing about it: whether a melody pleases is a matter of how it strikes the listener, and no argument can show that it ought to strike her otherwise. The second is that some judgments of taste are plainly wrong. A person who held that a forgotten versifier was the equal of Milton would not be expressing a harmless difference of sentiment; he would be mistaken, and the mistake would show something about his taste. Hume stated the tension with great clarity in \enquote{Of the Standard of Taste}. All sentiment is right, he observed, because sentiment refers to nothing beyond itself; and yet whoever would equate Ogilby with Milton, or Bunyan with Addison, would be thought to defend an extravagance as great as holding a molehill to be as high as Teneriffe \citep{Hume1757}.

Hume's resolution is that there is a standard of taste, a rule by which the various sentiments of men may be reconciled, and that it is fixed by the joint verdict of true judges: critics whose sentiment is delicate, whose judgment has been formed by practice and comparison, who are free of prejudice and who have the good sense to grasp what a work is and what it aims at. Such critics are rare, Hume conceded, and the question whether a given person is one is a question of fact, to be settled as other questions of fact are. But where they are found, their joint verdict is the standard, and a judgment of taste is correct in so far as it agrees with it.

The question we address is whether this standard is satisfiable: whether there is, for the judgments of taste on which criticism turns, a fact of the matter fixed by what ideal critics would jointly hold. AES-02 has argued that beauty is a response-dependent property, an object being beautiful just in case it is disposed to give aesthetic pleasure to qualified observers in suitable conditions, with the observers specified on Humean lines. That account settles the form of the facts about beauty, if there are any. It leaves open whether there are many such facts or few, since where qualified observers diverge the account yields only indeterminacy. The present manuscript closes that gap. We argue that the verdicts of ideal critics converge, and hence that there are facts about taste in every case on which the standard pronounces.

The argument has three parts. First, we state Hume's qualifications precisely, adding one that Hume's text implies but does not list separately: a moral sensitivity contained in good sense. Second, we classify the sources of divergence among judges and show that each is either removed by a qualification or confined to preferences between works of nearly equal merit. Third, we show that on these near-ties the joint verdict of ideal critics is a verdict of parity, so that the residual divergence is a divergence of preference and not of verdict. A further structural result bounds the disagreement of critics who meet the qualifications imperfectly, which explains why actual criticism, conducted by imperfect critics, approximates the standard rather than merely gesturing at it.

\section{The debate since Hutcheson}

The modern debate about the standard of taste is the debate about whether Hume's resolution works. We review the main proposals in order, with the consideration that has been thought to defeat each, since the argument of later sections is designed to meet those considerations.

\subsection{Internal sense and the eighteenth century}

\citet{Hutcheson1725} held that beauty is perceived by an internal sense, a natural power of receiving pleasure from uniformity amidst variety. Because the sense is natural and common to humankind, judgments of beauty are not arbitrary, and differences of taste are explained by the association of ideas, which can attach pleasure or disgust to objects for reasons that have nothing to do with their form. Hutcheson's account secures the commonality of taste at the price of its content: if the sense responds to uniformity amidst variety and to nothing else, it cannot explain why critics who agree that two works exhibit that property to the same degree still rank them very differently. \citet{Gerard1759} enlarged the account into a family of internal senses, of novelty, sublimity, beauty, imitation, harmony, ridicule and virtue, and made the refinement of taste a matter of their cultivation; but the multiplication of senses multiplied the possible sources of variation without saying which variations are errors. \citet{Dickie1996} gives the history of this period, and argues that Hume's essay is the strongest of its theories precisely because it does not tie taste to a single property.

Hume's essay shifts the question from the property to the judge. It does not say what beauty is in terms of the object; it says who is entitled to pronounce, and why. The five conditions, strong sense, delicate sentiment, practice, comparison and freedom from prejudice, are introduced as remedies for five ways in which a judge's sentiment can fail to be the sentiment the object is fitted to produce in a sound organ. The famous episode from \emph{Don Quixote}, in which Sancho's two kinsmen detect in a hogshead of wine the tastes of leather and of iron, later explained by a key on a leathern thong at the bottom, illustrates delicacy: a capacity to discriminate ingredients that are present in the object and that a coarser palate misses \citep{Hume1757}. Hume closes with an observation that will matter below. Two sources of variation, he says, remain even among good judges: the different humours of particular men, and the particular manners and opinions of an age and country. Where variation has these sources, there is no standard by which to reconcile it, and none is needed, since neither party is at fault.

\subsection{Kant and the claim to universality}

\citet{Kant1790} accepted the premise that the judgment of taste is grounded in a feeling of pleasure rather than in a concept of the object, and so cannot be proved. He nonetheless held that such a judgment claims the agreement of everyone, as a judgment of the merely agreeable does not, and he grounded the claim in the free play of imagination and understanding, faculties common to all who judge. The antinomy of taste, that judgments of taste are not based on concepts since otherwise they could be disputed by proofs, and are based on concepts since otherwise one could not even quarrel about them, is resolved by an indeterminate concept of the supersensible. Kant's account explains the normativity of the claim; it does not explain how the claim could be vindicated in a given dispute, since every actual disagreement can be attributed to a failure of one party to judge purely, and nothing in the account says which party it is. The Humean standard, by contrast, is in principle answerable to facts about judges, and that is why the later debate has been conducted mainly on Humean ground.

\subsection{The twentieth century: circles and double standards}

\citet{Sibley1959} argued that aesthetic concepts are not condition-governed: no description of an object in non-aesthetic terms entails that it is graceful or garish, although aesthetic properties depend on non-aesthetic ones. The application of aesthetic concepts requires taste, a capacity that goes beyond normal eyesight and hearing. Sibley's thesis is often read as supporting the Humean picture, since it locates aesthetic competence in a special sensitivity, and as undermining it, since it denies that the critic can prove her verdicts. \citet{Mothersill1984} drew the second conclusion more sharply: there are no laws of taste, no principles from which verdicts follow, and yet judgments of taste are genuine judgments with truth values. Her view leaves the existence of facts about taste and the possibility of a standard to be argued for separately.

The first systematic objection to Hume's argument concerns circularity. If true judges are identified by the soundness of their verdicts, and sound verdicts by their agreement with true judges, the standard is empty. \citet{Kivy1967} argued that the circle can be broken, because the qualifications can be specified and tested independently of the verdicts: delicacy, in particular, is testable by the discrimination of ingredients, as in the case of the hogshead. \citet{Carroll1984} pressed the objection in another form, arguing that Hume's account is best read as a causal theory of which objects produce pleasure in normal observers, and that so read it cannot deliver the normative standard Hume wanted, since the critic's verdict is evidence of the causal power of the object and not a ground of correctness. \citet{Shelley1994} argued that Hume operates with a double standard: the joint verdict of true judges and the test of time, the durable admiration that works such as Homer's command across ages and nations, and that the two can come apart. AES-02 has given the reply to the circularity objection that we shall rely on: the qualifications are specifiable without the concept of beauty, and the fact that critics are in practice identified by their verdicts is an epistemic fact about how we find them, not a constitutive fact about what makes them qualified.

\citet{Levinson2002} formulated what he called the real problem with Hume's standard. Suppose that the joint verdict of true judges ranks one work above another. Why should an ordinary person, whose own preferences run the other way and who is not a true judge, care about the verdict? The standard may be fixed, but it is not clear what reason it gives to anyone who does not already share it. Levinson's answer was that true judges are better guides to the experiences that are most rewarding, so that the ordinary person has reason to follow them in so far as she wants such experiences; his worry was that the answer seems to presuppose that the rewards are ones the ordinary person can come to have. We take up the problem in Section~\ref{s:obj}.

\subsection{The relativist challenge}

The most serious objection to the satisfiability of the standard concerns convergence. \citet{Goldman1995} argued that ideal critics, defined by Hume's qualifications, will diverge in their evaluations, because their tastes have been formed by different exposures and reflect different temperaments, and that aesthetic evaluations are therefore true only relative to a sensibility. \citet{Bender1996} argued that disagreements among competent critics about the aesthetic properties of works are in some cases irresolvable, and that the best explanation of irresolvability is that there is no single fact to be right about; he took this to tell against aesthetic realism, while allowing that aesthetic properties supervene on non-aesthetic ones relative to a sensibility. \citet{Levinson2001} argued that ideal critics may differ in sensibility, so that the evaluative force of some aesthetic properties, whether elegance or grandeur counts in a work's favour and how much, varies between them. On these views Hume's own concession about humours and manners is not a residue but the main phenomenon: the standard is satisfiable only within communities of sensibility, and there are facts about taste only relative to such communities.

The relativist argument has two premises: that differences of sensibility survive idealization, and that such differences produce divergent verdicts on the same works. We grant the first premise and deny the second. Differences of sensibility survive idealization, as Hume said; but they produce divergent \emph{preferences}, not divergent \emph{verdicts}, except where the merits of the works compared are close enough that an ideal critic does not issue a verdict of superiority at all.

\subsection{Recent literature}

Work of the last fifteen years has moved in three directions. The first concerns the relation between real and ideal critics. \citet{Ross2008} argued that Hume's true judges should be read as real people, identifiable in practice, rather than as idealizations, and that the standard is accordingly a matter of the actual joint verdict of the best available critics. \citet{Kieran2008} argued that ideal critics, as usually described, are not ideal: the qualifications concern discrimination and knowledge, but the value of aesthetic experience also depends on the critic's motivations and character, which the Humean list does not constrain. \citet{Shelley2013} examined what the joint verdict of true judges is, arguing that it is not a summary of what each judge says but the verdict on which judges would agree. Our account takes a position on each of these questions: the standard is fixed by ideal critics, actual critics approximate it to a degree that can be bounded (Section~\ref{s:bound}), and good sense includes moral sensitivity, which answers part of Kieran's worry.

The second direction concerns aesthetic testimony. \citet{Hopkins2011} defended a pessimism on which aesthetic judgments formed on testimony are illegitimate, even when the testifier is reliable, and argued that the best explanation of this norm is not that there are no aesthetic facts to transmit but that aesthetic judgment is subject to a requirement of autonomy. \citet{Robson2012} surveys the debate between pessimists and optimists. The relevance to the standard of taste is direct: if the joint verdict of ideal critics fixes facts about taste, it would seem that one could come to know those facts by deferring to the critics, and the apparent illegitimacy of deference has been taken as evidence that there are no such facts.

The third direction is empirical. \citet{Cutting2003} found that the canon of French Impressionism correlates with the frequency with which works have been reproduced, and that mere exposure can shift the preferences of naive viewers among comparable paintings; he suggested that canon formation is driven in part by exposure. \citet{Meskin2013} tested the hypothesis with works of very different merit, paintings by Millais and by Kinkade, and found that repeated exposure increased liking for the former and decreased liking for the latter. \citet{Vessel2018} found that agreement among individuals in aesthetic judgments is much higher for natural objects such as faces and landscapes than for artefacts such as architecture and abstract art. Each of these findings bears on whether convergence is to be expected, and we return to them in Sections~\ref{s:cases} and~\ref{s:obj}.

\section{Framework}\label{s:frame}

We distinguish two kinds of evaluative response that a judge may make to works of art and other objects of taste.

\begin{definition}[Verdict]\label{d:verdict}
A \emph{verdict} is a judgment of aesthetic merit made with a claim to correctness: that a work is beautiful, or excellent of its kind, or that one work is aesthetically better than another. An \emph{absolute} verdict concerns one work; a \emph{comparative} verdict concerns two.
\end{definition}

\begin{definition}[Preference]\label{d:pref}
A \emph{preference} is a disposition to take greater pleasure in, or to seek out, one work rather than another, avowed without a claim that others ought to share it.
\end{definition}

The distinction is familiar from ordinary critical practice. A critic may hold that Bach's cantatas and Handel's oratorios are both supreme achievements of their kind, and that neither is better than the other, while acknowledging that she would rather listen to Handel. Hume draws the distinction himself when he notes that at twenty one may prefer Ovid, at forty Horace and at fifty Tacitus, and that it would be vain to dispute about such preferences, since each man is pleased by what answers to his disposition \citep{Hume1757}. What Hume does not say, and what the argument below requires, is that the true judge keeps the two apart.

\begin{definition}[Ideal critic]\label{d:ideal}
An \emph{ideal critic} is a judge who satisfies the following conditions:
\begin{enumerate}[label=(\roman*),nosep]
\item \emph{Delicacy}: she discriminates all the non-aesthetic features of a work relevant to its effect, at the finest grain at which they make a difference to the response of a sound organ;
\item \emph{Practice}: she has attended repeatedly to works of the relevant kind, so that her response to a work is not distorted by novelty or by the first confusion of a complex object;
\item \emph{Comparison}: she has acquaintance with a wide range of works in the relevant categories and in neighbouring ones, so that she can place each work in its correct category and estimate its merit relative to others;
\item \emph{Freedom from prejudice}: she has no interest in the work's author, owner, reputation or consequences that bears on her response, and she places herself in the point of view of the audience the work supposes;
\item \emph{Good sense}: she understands the work's ends, its conventions and its representational content, including the moral perspective it adopts and the responses it prescribes, and she can assess whether those responses are merited; and she discriminates between her verdicts and her preferences, issuing a verdict only where her assessment does not depend on her particular temperament.
\end{enumerate}
\end{definition}

Conditions (i) to (iv) and the first half of (v) are Hume's \citep{Hume1757}. Hume attributes to good sense the capacity to grasp the relations of the parts of a work, its end, and the consistency of its characters; he also holds that where a work represents vicious manners without marking them with disapproval, this is a real deformity, which the critic must not excuse by appeal to the manners of the age. A critic who could not see that a work's moral perspective is unacceptable would therefore fail to detect a deformity that Hume counts as one. We make this explicit: good sense includes moral sensitivity, the capacity to recognize the moral perspective a work adopts and whether the responses it prescribes are merited. This is the sense of \enquote{ideal critic} that AES-08 relies on when it holds that a moral flaw is an aesthetic flaw when it prevents the response a work prescribes from a sensitive and informed audience. The second half of (v), the discrimination between verdict and preference, is the condition we add. It is a form of self-knowledge acquired by comparison: a critic learns how her temperament inflects her responses by comparing them with those of other practiced judges whose temperaments differ, as a musician learns the bias of her own ear.

\begin{definition}[Joint verdict and convergence]\label{d:conv}
The verdicts of ideal critics \emph{converge} on a question of taste if every ideal critic who addresses the question in suitable conditions issues the same verdict on it. Where they converge, their common verdict is the \emph{joint verdict}. Hume's standard of taste is \emph{satisfiable} if (a) the conditions of Definition~\ref{d:ideal} are specifiable without the concepts of beauty or aesthetic merit, (b) they can be approximated by actual judges to degrees that make a difference to the reliability of their verdicts, and (c) the verdicts of ideal critics converge on every question of taste they address.
\end{definition}

Three principles state the background of the argument.

\begin{principle}[Common nature]\label{p:common}
Aesthetic responses are produced by a sensibility common to normal human beings: the same features of an object, perceived in the same category and the same conditions, are disposed to produce the same aesthetic response in all observers whose perceptual and affective capacities are normal, except as modified by the sources listed in Principle~\ref{p:sources}.
\end{principle}

\begin{principle}[Sources of divergence]\label{p:sources}
Every divergence in aesthetic response between two observers attending to the same work arises from one or more of the following: (a) differences in the discrimination of the work's features; (b) differences in familiarity with the work and works of its kind; (c) differences in the category in which the work is perceived, or in knowledge of its conventions; (d) interests, prejudices or associations bearing on the work; (e) differences in moral perspective or in understanding of the work's content; (f) differences of temperament, including those of age and character.
\end{principle}

\begin{principle}[Categories]\label{p:cat}
The aesthetic merit of a work is its merit as perceived in its correct category, and the correct category is fixed by facts about the work's origin and the practices in which it was made \citep{Walton1970}.
\end{principle}

Principle~\ref{p:common} is Hume's claim that the general principles of taste are uniform in human nature. It is an empirical premise, and we do not claim to establish it a priori; we note in Section~\ref{s:cases} what the evidence for it is, and in Section~\ref{s:limits} what follows if it fails. Principle~\ref{p:sources} expands Hume's list of defects with the manners of an age divided into its two components: unfamiliarity with the conventions of another time or place, which is category (c), and differences of moral opinion, which is part of (e). The list is offered as exhaustive. Any proposed source of divergence not on it must either be a variety of one of these, or be a source that operates independently of the work's features and the observer's capacities, in which case it is not a source of divergence in aesthetic response to the work. Principle~\ref{p:cat} is the thesis of AES-04 and AES-16 that aesthetic value is assessed in a work's correct category, and of AES-11 that the categories of natural objects are fixed by natural history.

The influence of temperament is limited. We state the limit as an assumption, which Hume's examples support and which we defend in Section~\ref{s:cases}.

\begin{assumption}[Bounded temperament]\label{a:temp}
Differences of temperament between observers who are otherwise ideal alter the weight each gives to particular merits, such as tenderness or severity, brilliance or gravity, but they do not reverse the ranking of two works whose merit differs by more than a bounded amount. Call this bound the \emph{margin of temperament}.
\end{assumption}

Hume's illustrations are all of this form. The young man and the old man differ in whether they prefer Ovid or Tacitus; neither prefers Ogilby to Milton. One person likes the sublime and another the tender; neither takes a sentimental ballad to be greater than a Schubert song because of it. Temperament redistributes pleasure among works that offer different excellences in comparable degree, and Assumption~\ref{a:temp} says only that there is a degree of difference beyond which it does not do so.

\section{The argument}\label{s:arg}

\begin{lemma}[Non-circularity]\label{l:noncirc}
The conditions of Definition~\ref{d:ideal} can be specified without the concepts of beauty or aesthetic merit.
\end{lemma}

\begin{proof}
Delicacy is a capacity to discriminate non-aesthetic features: pitches, intervals, hues, proportions, rhythms, the order of events in a narrative. Practice and comparison are matters of exposure and can be specified historically. Freedom from prejudice is the absence of certain interests and associations, and the adoption of the point of view of a work's intended audience is specified by the facts about that audience. Good sense is understanding of a work's ends, conventions and content, which are fixed by its origin and its features, together with moral sensitivity, which is specified by moral concepts and not aesthetic ones; the discrimination of verdict from preference is a matter of self-knowledge of temperament, acquired by comparison. None of these specifications uses the concept of beauty or of aesthetic merit. This is the argument of AES-02, extended to the fifth condition.
\end{proof}

\begin{lemma}[Eliminable sources]\label{l:elim}
Two ideal critics do not diverge in aesthetic response to a work from sources (a) to (d) of Principle~\ref{p:sources}.
\end{lemma}

\begin{proof}
Source (a) is removed by delicacy, since each critic discriminates every feature relevant to the response of a sound organ. Source (b) is removed by practice: each has attended to the work and to works of its kind sufficiently to have overcome the distortions of first acquaintance. Source (c) is removed by comparison and by the requirement to place oneself in the point of view of the intended audience: each perceives the work in its correct category, which by Principle~\ref{p:cat} is fixed by facts that both know, and each knows the conventions under which the work was made. Source (d) is removed by freedom from prejudice. Since both critics meet all four conditions, no divergence between them can arise from these sources.
\end{proof}

\begin{lemma}[Moral perspective]\label{l:moral}
Two ideal critics do not diverge in aesthetic verdict from source (e) of Principle~\ref{p:sources}, except in so far as there is no fact about the moral question at issue.
\end{lemma}

\begin{proof}
The part of (e) that concerns understanding of the work's content is removed by good sense. The remainder concerns moral perspective. By the account of AES-08, a work's moral perspective bears on its aesthetic merit when, and only when, it prevents a sensitive audience from giving the response the work prescribes. Whether it does so depends on whether the prescribed response is merited, which is a moral question. Ideal critics are morally sensitive (Definition~\ref{d:ideal}(v)), so where there is a fact about whether the response is merited, each recognizes it, and their verdicts on whether the work fails on its own terms agree. A divergence can therefore persist only where there is no fact about the moral question, and nothing in the argument requires that there be such questions bearing on works of art. MTE-12 argues that moral disagreement is no deeper than disagreement in other theoretical domains that are not thought to lack facts; if so, the exception is no wider than the exception for indeterminate questions in any such domain.
\end{proof}

The remaining source is temperament, which the qualifications do not remove and which Hume rightly held cannot be removed. The next lemma shows that it does not produce divergent verdicts.

\begin{lemma}[Temperament]\label{l:temp}
Differences of temperament between ideal critics produce divergent preferences, but not divergent verdicts.
\end{lemma}

\begin{proof}
Let two ideal critics address a comparative question between works $x$ and $y$, and suppose that, by Lemmas~\ref{l:elim} and~\ref{l:moral}, their responses differ only by the influence of temperament. There are two cases. If the merit of $x$ and of $y$ differs by more than the margin of temperament, then by Assumption~\ref{a:temp} temperament does not reverse the ranking, and both critics rank the works the same way. If the difference is within the margin, then each critic, by the second half of Definition~\ref{d:ideal}(v), knows that her ranking of the two works depends on her temperament, and so does not issue a verdict of superiority; she judges the works of comparable merit, and avows her ranking as a preference. Both critics then issue the verdict of parity, while their preferences may differ. The argument for absolute verdicts is the same, with the threshold for the verdict in place of the second work: where a work's merit lies within the margin of temperament of the threshold, an ideal critic does not issue the verdict that it meets or fails to meet it, and records instead that it lies at the boundary. In no case do the verdicts differ.
\end{proof}

\begin{theorem}[Convergence]\label{t:main}
The verdicts of ideal critics converge on every question of taste they address.
\end{theorem}

\begin{proof}
Suppose two ideal critics address the same question about the same work or works in suitable conditions. By Principle~\ref{p:common}, their responses can differ only through the sources listed in Principle~\ref{p:sources}. By Lemma~\ref{l:elim}, sources (a) to (d) produce no divergence; by Lemma~\ref{l:moral}, source (e) produces no divergence of verdict on any question where there is a fact for their moral sensitivity to track; by Lemma~\ref{l:temp}, source (f) produces divergent preferences but not divergent verdicts. Hence their verdicts agree. Since the two critics were arbitrary, every ideal critic who addresses the question issues the same verdict, which by Definition~\ref{d:conv} is convergence.
\end{proof}

\begin{corollary}[Satisfiability]\label{c:satisf}
Hume's standard of taste is satisfiable.
\end{corollary}

\begin{proof}
Condition (a) of Definition~\ref{d:conv} is Lemma~\ref{l:noncirc}, and condition (c) is Theorem~\ref{t:main}. Condition (b), approximability, is established by Proposition~\ref{p:approx} below.
\end{proof}

\begin{corollary}[Facts about taste]\label{c:facts}
There are facts about taste. For every question of taste addressed by ideal critics, there is a correct answer, namely their joint verdict, and an ordinary judgment of taste is correct just in case it agrees with that verdict.
\end{corollary}

\begin{proof}
By Theorem~\ref{t:main}, the joint verdict exists for every question ideal critics address. On the response-dependent account of AES-02, whether an object is beautiful is fixed by its disposition to please qualified observers in suitable conditions, and by Lemma~\ref{l:temp} that disposition is determinate wherever ideal critics issue an absolute verdict: they agree that the object meets the threshold, or that it fails to meet it, or that it lies at the boundary. Comparative facts are fixed in the same way. A judgment that agrees with the joint verdict is therefore true, and one that disagrees with it is false.
\end{proof}

Corollary~\ref{c:facts} includes among the facts about taste the facts of parity. That Ovid and Tacitus are each excellent in their kind and that neither is better is not an absence of fact but a fact, on which ideal critics agree; what they do not agree on is which they would rather read. This is how the result is consistent with AES-02's treatment of divergence. AES-02 held that where qualified observers diverge about whether an object is beautiful, the question has no determinate answer. Theorem~\ref{t:main} shows that ideal critics do not diverge in verdict; the indeterminacy AES-02 envisaged is located at the boundary of the threshold, where the joint verdict is precisely that the work is a borderline case.

\section{A bound on disagreement}\label{s:bound}

The convergence theorem concerns ideal critics, and no actual critic is ideal. If the standard were only a fact about idealized judges, with no bearing on the judgments of actual ones, the objection of \citet{Ross2008} would have force: a standard that no one can apply is not a standard of the practice of criticism. We now show that actual critics approximate the standard, in a sense that can be made precise.

\begin{assumption}[Graded deficiency]\label{a:graded}
The aesthetic merit of works in a category can be represented on an interval scale, and each observer $o$ has a degree of deficiency $d_k(o)\ge 0$ in each of the conditions $k$ of Definition~\ref{d:ideal}, with $d_k(o)=0$ for all $k$ just in case $o$ is ideal. There are constants $L>0$ and $\tau\ge 0$ such that for every work $x$ the response of $o$ to $x$ differs from the merit $m(x)$ of $x$ by at most $L\sum_k d_k(o)+\tau$, where $\tau$ is the margin of temperament.
\end{assumption}

Assumption~\ref{a:graded} adds to Principle~\ref{p:sources} a claim of proportionality: the less a judge falls short of each condition, the less her response can depart from the merit of the work on account of that condition. The interval scale is an idealization; the argument uses only the inequalities, and would go through with any representation on which differences of merit can be compared with the bounds.

\begin{proposition}[Agreement of imperfect critics]\label{p:approx}
Let $\delta\ge 0$, and let $x$ and $y$ be works in the same category with $m(x)-m(y)>2(L\delta+\tau)$. Every observer whose total deficiency $\sum_k d_k(o)$ is at most $\delta$ responds more favourably to $x$ than to $y$. In particular, the set of comparative questions on which all observers of total deficiency at most $\delta$ agree with the joint verdict increases as $\delta$ decreases, and in the limit $\delta=0$ it contains every pair whose merit differs by more than twice the margin of temperament.
\end{proposition}

\begin{proof}
Write $r_o(x)$ for the response of $o$ to $x$. By Assumption~\ref{a:graded}, $r_o(x)\ge m(x)-L\delta-\tau$ and $r_o(y)\le m(y)+L\delta+\tau$. Hence $r_o(x)-r_o(y)\ge m(x)-m(y)-2(L\delta+\tau)>0$. The set of pairs satisfying the hypothesis grows as $\delta$ decreases, since the bound $2(L\delta+\tau)$ decreases. At $\delta=0$ the bound is $2\tau$, and on every pair whose merit differs by more than that, ideal critics rank the works as merit does, which by Lemma~\ref{l:temp} is the joint verdict.
\end{proof}

\begin{corollary}[Approximation]\label{c:approx}
Actual critics approximate the standard of taste: the less they fall short of Hume's conditions, the larger the class of questions on which their verdicts are guaranteed to agree with the joint verdict and with one another. Disagreement among critics of small total deficiency is confined to pairs of works whose merit differs by little.
\end{corollary}

Corollary~\ref{c:approx} is condition (b) of Definition~\ref{d:conv}, and it explains a feature of critical practice that a skeptic about the standard must explain otherwise. Disputes among well-qualified critics are rarely about whether Milton is better than Ogilby. They are about whether a particular late quartet is the equal of another, whether a novelist's later manner is an advance on her earlier one, whether a new building succeeds. Proposition~\ref{p:approx} predicts exactly this distribution: as qualification increases, disagreement retreats to close cases. It also explains the authority of critics without making it absolute. A critic's verdict on a question with a wide gap of merit is very likely correct even if she is far from ideal; her verdict on a close question carries little weight unless her deficiency is small. The ordinary practice of trusting critics on large questions and arguing with them on small ones is the practice the proposition recommends.

The proposition also bears on the double standard of \citet{Shelley1994}. The test of time is the agreement of many judges, of different periods and nations, each imperfect in different ways. Their deficiencies differ, and the errors they induce, of fashion, of local prejudice, of unfamiliarity with older conventions, are not systematically aligned. A work that wins the agreement of judges whose deficiencies are differently distributed is a work whose merit exceeds the bounds those deficiencies impose. The test of time is therefore not a second standard but evidence about the first: it identifies works whose merit is high enough relative to their competitors that imperfect judges of every period agree in ranking them highly. That Homer, as Hume put it, still pleases at Paris and at London as he pleased at Athens and Rome is evidence of the merit that ideal critics would recognize, and is explained by it.

\section{Further cases}\label{s:cases}

We now examine a series of cases, each of which has been taken to show that the standard fails, and analyze each with the principles of Section~\ref{s:frame}.

\paragraph{Ogilby and Milton.} Hume's paradigm is a comparison across a very large gap of merit. Every qualified reader of English verse ranks \emph{Paradise Lost} above Ogilby's translations, and the ranking is stable across centuries and across readers of very different temperaments and opinions. The case falls under the first clause of Lemma~\ref{l:temp}: no difference of temperament reverses so great a difference, and the observers who would reverse it are those deficient in practice and comparison. By Proposition~\ref{p:approx}, the verdict is shared even by quite imperfect judges.

\paragraph{Ovid, Horace and Tacitus.} Hume's own example of blameless variation is a comparison across kinds, between a poet of love, a poet of measured reflection and a historian of political darkness. Each is excellent, and their excellences are different. The case falls under the second clause of Lemma~\ref{l:temp}: an ideal critic does not hold that Tacitus is better than Ovid, but that both are among the masters of Latin prose and verse, while she may, at fifty, prefer to read Tacitus. The relativist describes this case as a divergence of verdict between equally ideal judges. It is a divergence of preference between judges who agree in verdict, and critics do in fact describe such cases in that way.

\paragraph{The first reception of the new.} The premiere of \emph{The Rite of Spring} in Paris in 1913 provoked an uproar, and many of the first reviews were hostile; within a generation the work was regarded as one of the central achievements of twentieth-century music. The case is sometimes cited to show that critics of a period, including well-qualified ones, can be radically wrong, and the inference is drawn that there is nothing to be right about. Our principles explain the case without that inference. The first audience lacked practice with the work and had no comparison class for it; its category, a new idiom of rhythm and orchestration, was not one in which they had learned to hear. The divergence was from sources (b) and (c), which are eliminable, and it was eliminated as practice and comparison accumulated. The case is evidence that actual critics are imperfect, which no one doubts, and that their convergence improves with the qualifications Hume lists, which is what Proposition~\ref{p:approx} predicts.

\paragraph{Neglect and revival.} The reputation of Vermeer was slight for nearly two centuries after his death, until the critic Th\'eophile Thor\'e-B\"urger championed him in the 1860s; Bach's \emph{St Matthew Passion} was rarely performed after his death until Mendelssohn's Berlin performance of 1829. Revivals of this kind are revisions against prevailing prestige, and they are conducted by appeal to features of the works that earlier judges had neglected: Vermeer's handling of light, Bach's architecture of chorus and recitative. They are corrections of source (a) and of source (d), the last in the form of a prejudice against what was unfashionable. They show that the standard can come apart from the actual consensus of a period, which is a condition of its being a standard rather than a report of consensus.

\paragraph{Unfamiliar traditions.} A listener trained in Western tonal music who first hears the court music of another tradition may find it monotonous or formless, and a listener trained in that tradition may find a Romantic symphony bombastic. If such reactions were verdicts of ideal critics, they would refute convergence. They are not, because they fail condition (iii): each listener lacks comparison in the other's category, and so hears the music in a category that is not its correct one. Principle~\ref{p:cat} requires that a work be assessed in its correct category, and an ideal critic is one who can perceive it so. Where listeners have acquired competence in both traditions, the evidence of ethnomusicology and of practicing musicians trained in both is that they recognize excellence in each by the standards of each, and do not regard the traditions as competing for a single ranking. Comparisons across categories are then of the kind discussed under Ovid and Tacitus.

\paragraph{The morally flawed masterpiece.} A film made to celebrate a murderous regime may be formally brilliant. A critic insensitive to its moral perspective will rank it very highly; a morally sensitive critic will judge that it fails as a work of celebration, since the admiration it prescribes is not merited, while acknowledging its formal achievement. If moral sensitivity were not among the qualifications, ideal critics might diverge on such works according to their moral opinions, and some would be right only by luck. Definition~\ref{d:ideal}(v) excludes this. The case shows why the condition belongs on the list: without it, critics who meet the other four conditions would diverge in verdict on a class of works that the practice of criticism does not regard as indeterminate.

\paragraph{Exposure and kitsch.} The finding of \citet{Cutting2003}, that exposure shifts the preferences of naive viewers among comparable Impressionist paintings, is what Assumption~\ref{a:graded} predicts for observers deficient in practice and comparison considering works whose merit differs little: for such pairs, the bound of Proposition~\ref{p:approx} is not met, and responses are free to vary with factors such as familiarity. The finding of \citet{Meskin2013}, that exposure increases liking for Millais and decreases it for Kinkade, is what Hume's condition of practice predicts for works of very different merit: repeated attention reveals the merits of the better work and the defects of the worse, so that practice brings judges toward the joint verdict. Together the two findings suggest that exposure is not a rival explanation of convergence on the canon, but a mechanism through which practice operates.

\paragraph{Faces and artefacts.} \citet{Vessel2018} report that agreement among untrained individuals is much higher for faces and landscapes than for artworks and buildings. The finding is consistent with Principle~\ref{p:common}, and we take it to support it. For natural objects such as faces, nearly everyone has extensive practice and comparison, acquired in ordinary life; for artworks, the qualifications require training that most of the participants lacked. The prediction of the account is that agreement in the second domain rises with training in the relevant categories, and it is that prediction, rather than the level of agreement among untrained observers, that bears on the satisfiability of the standard.

\section{Objections and replies}\label{s:obj}

\begin{objection}[Idealization]\label{o:ideal}
No actual critic satisfies Definition~\ref{d:ideal}, and perhaps none could. A standard fixed by ideal critics is then a standard no one can know or apply, and the claim that there are facts about taste is idle \citep{Ross2008}.
\end{objection}

\begin{reply}
The standard is fixed by ideal critics, but it is known through actual ones. Corollary~\ref{c:approx} shows that actual critics' verdicts agree with the joint verdict on a class of questions that grows as their deficiency shrinks; we can therefore know the standard on many questions by consulting critics who are far from ideal, and on closer questions by consulting better ones and by attending to the reasons they give. The situation is that of any measurement by imperfect instruments: no thermometer is perfect, and yet temperatures are known, to a precision that depends on the instrument. Ross's proposal that the standard be the actual verdict of the best available critics would make the best critics of each period infallible about their own period's works, and so could not explain the corrections illustrated by the revivals of Vermeer and Bach.
\end{reply}

\begin{objection}[Ideal critics are not ideal]\label{o:kieran}
Hume's qualifications concern discrimination and knowledge. But the value of a work also depends on how it engages a person's character and motivations, and critics who are equally discriminating may differ in these respects in ways that bear on value. The idealized critic is a connoisseur, not a judge of what matters in art \citep{Kieran2008}.
\end{objection}

\begin{reply}
The objection identifies a real gap in the usual reading of Hume, and our account closes part of it. Good sense, as Definition~\ref{d:ideal}(v) states it, includes moral sensitivity and an understanding of the responses a work prescribes; an ideal critic is therefore not merely discriminating but responsive to what is at stake in a work. The rest of the gap concerns differences of character and motivation that are not moral, such as a disposition to value austerity over abundance. These are differences of temperament, and by Lemma~\ref{l:temp} they bear on preferences and on close comparisons, not on verdicts across large differences of merit. The objection shows that ideal critics need more than a fine palate; it does not show that critics who have what they need will diverge.
\end{reply}

\begin{objection}[Divergent sensibilities]\label{o:rel}
Ideal critics may differ in sensibility in ways that alter the evaluative force of aesthetic properties: for one, grandeur weighs more than delicacy; for another, the reverse. Their verdicts on works that combine these properties in different proportions will then diverge, and the divergence is blameless. Aesthetic judgments are true only relative to a sensibility \citep{Goldman1995,Bender1996,Levinson2001}.
\end{objection}

\begin{reply}
The objection describes differences of temperament, and we have granted that they survive idealization. What it must show is that they produce divergent verdicts on works whose merit differs by more than the margin of temperament, and its examples do not show this. The cases offered are comparisons between works that excel in different respects to comparable degrees, of grandeur against delicacy, of the severe against the lyrical; these are the cases in which Lemma~\ref{l:temp} predicts parity of verdict and divergence of preference. A relativist who holds that the margin of temperament is so large that it covers all or most comparisons must deny Assumption~\ref{a:temp} in a strong form, and then must explain the very wide agreement among qualified critics, across temperaments, on large comparisons within and across categories. The test of time and the distribution of disagreement among critics noted after Corollary~\ref{c:approx} are both evidence that the margin is small relative to the range of merit in each art.
\end{reply}

\begin{objection}[Sociology of taste]\label{o:bourdieu}
The convergence of critics is a social fact. Taste functions to mark distinctions between classes, and the agreement of critics reflects their membership of a class whose taste is legitimated as the standard. The qualifications of the true judge are the competences of the dominant fraction, and their joint verdict is the verdict of that fraction \citep{Bourdieu1984}.
\end{objection}

\begin{reply}
That the acquisition of the qualifications is unequally distributed, and that taste is used to mark social distinction, are true and important. Neither shows that the verdicts of those who have the qualifications are fixed by social position rather than by the works. A sociological explanation of convergence predicts that critics agree on whatever the dominant group favours, and that their verdicts change with the composition of that group. It does not predict the cross-period stability of the test of time, since the dominant groups of Athens and of eighteenth-century London had little in common; nor the revivals of neglected work against prevailing prestige; nor the finding of \citet{Meskin2013} that practice moves untrained observers toward the critical consensus on works of very different merit. The objection is best taken as an argument for wider access to practice and comparison, which the account supports, since by Corollary~\ref{c:approx} wider access widens the class of people whose verdicts agree with the standard.
\end{reply}

\begin{objection}[The real problem]\label{o:levinson}
Even if the joint verdict of ideal critics exists, an ordinary person whose preferences run against it has no reason to defer to it. The standard is a fact about a small class of judges, not a norm for anyone else \citep{Levinson2002}.
\end{objection}

\begin{reply}
The standard does not require anyone to change her preferences, and it does not require deference. It fixes which verdicts are true. The ordinary person who judges a sentimental ballad better than a Schubert song is mistaken, by Corollary~\ref{c:facts}, whatever she prefers; she is not thereby obliged to prefer the song. The reason the standard gives her is the reason any truth gives: if she cares about judging correctly, she has reason to judge in accordance with it. She also has a further reason, which is Levinson's: the works the joint verdict ranks highest are those that offer the greatest aesthetic rewards to anyone who acquires the qualifications, and by Corollary~\ref{c:approx} those qualifications can be approached by degrees. Whether she wants those rewards is up to her. That the standard does not compel desire is not a defect in a standard of correctness.
\end{reply}

The question of aesthetic testimony, raised in Section~2, is answered by the same distinction. If the joint verdict fixes facts, one can come to know some of them by testimony, as EPI-08 holds that testimony transmits knowledge in general. What testimony cannot transmit is the appreciation of a work, which requires the response itself; and a judgment of taste, as distinct from a belief about the joint verdict, is a judgment made from one's own response. The norm against judging on testimony that \citet{Hopkins2011} describes is therefore a norm on judgments of taste, not evidence that there is nothing for testimony to transmit.

\section{What the result does not establish}\label{s:limits}

The result does not provide a decision procedure for questions of taste. It shows that there is a joint verdict on every question ideal critics address, and that actual critics approximate it; it does not tell anyone how to settle a close question, and on such questions the reasons critics give, pointing to features of the work and comparing it with others, remain the only method.

It does not establish Principle~\ref{p:common} or Assumption~\ref{a:temp} a priori. Both are empirical, and the argument is as strong as the evidence for them: the test of time, the distribution of disagreement among critics, and the effects of practice on the judgments of untrained observers. If human sensibility were divided into kinds that differ in their response to the same features beyond any bounded margin, the convergence theorem would hold within each kind and fail across them, and the facts about taste would be relative to kind. The evidence reviewed in Section~\ref{s:cases} gives no reason to think that this is our situation.

It does not make actual critics authorities. Their verdicts are evidence of the joint verdict, of a strength that depends on their qualifications and on the size of the difference of merit at issue. It does not show that preferences are subject to a standard: where works are of comparable merit, the standard is silent about which a person should enjoy more, and Hume's remark that it is vain to dispute about such preferences stands. Nor does it extend to the merely agreeable. The sources of divergence listed in Principle~\ref{p:sources} concern aesthetic responses to works and objects perceived in their categories; the pleasure of a food or a scent may vary with constitution in ways that no qualification removes, and the argument says nothing about it.

Finally, the account is consistent with the claim of AES-16 that there is no progress in artistic value. The standard ranks works by their merit in their categories and as achievements in their situations, and nothing in it implies that later works rank higher.

\section{Conclusion}

Hume's standard of taste is satisfiable. The qualifications of the true judge can be specified without the concepts of beauty or merit, actual critics approximate them to degrees that make a measurable difference to the reliability of their verdicts, and the verdicts of ideal critics converge. Convergence follows once two points are made explicit: that good sense includes the moral sensitivity needed to judge whether the responses a work prescribes are merited, and that the true judge distinguishes her verdicts from her preferences. Every source of divergence among judges is then either removed by a qualification, or confined, in the case of temperament, to preferences among works of nearly equal merit, on which the joint verdict is parity. There are therefore facts about taste, including facts of parity, and an ordinary judgment of taste is correct when it agrees with the joint verdict that these facts constitute.

\begin{thebibliography}{99}\small
\bibitem[Bender(1996)]{Bender1996} Bender, J.~W. (1996). Realism, supervenience, and irresolvable aesthetic disagreements. \emph{The Journal of Aesthetics and Art Criticism}, 54(4), 371--381.
\bibitem[Bourdieu(1984)]{Bourdieu1984} Bourdieu, P. (1984). \emph{Distinction: A Social Critique of the Judgement of Taste}. Trans. R. Nice. Harvard University Press. (Original work published 1979.)
\bibitem[Carroll(1984)]{Carroll1984} Carroll, N. (1984). Hume's standard of taste. \emph{The Journal of Aesthetics and Art Criticism}, 43(2), 181--194.
\bibitem[Cutting(2003)]{Cutting2003} Cutting, J.~E. (2003). Gustave Caillebotte, French Impressionism, and mere exposure. \emph{Psychonomic Bulletin \& Review}, 10(2), 319--343.
\bibitem[Dickie(1996)]{Dickie1996} Dickie, G. (1996). \emph{The Century of Taste: The Philosophical Odyssey of Taste in the Eighteenth Century}. Oxford University Press.
\bibitem[Gerard(1759)]{Gerard1759} Gerard, A. (1759). \emph{An Essay on Taste}. A. Millar.
\bibitem[Goldman(1995)]{Goldman1995} Goldman, A.~H. (1995). \emph{Aesthetic Value}. Westview Press.
\bibitem[Hopkins(2011)]{Hopkins2011} Hopkins, R. (2011). How to be a pessimist about aesthetic testimony. \emph{The Journal of Philosophy}, 108(3), 138--157.
\bibitem[Hume(1757)]{Hume1757} Hume, D. (1757). Of the standard of taste. In \emph{Four Dissertations}. A. Millar.
\bibitem[Hutcheson(1725)]{Hutcheson1725} Hutcheson, F. (1725). \emph{An Inquiry into the Original of Our Ideas of Beauty and Virtue}. J. Darby.
\bibitem[Kant(1790)]{Kant1790} Kant, I. (1790). \emph{Critique of the Power of Judgment}. Trans. P. Guyer and E. Matthews. Cambridge University Press, 2000.
\bibitem[Kieran(2008)]{Kieran2008} Kieran, M. (2008). Why ideal critics are not ideal: Aesthetic character, motivation and value. \emph{The British Journal of Aesthetics}, 48(3), 278--294.
\bibitem[Kivy(1967)]{Kivy1967} Kivy, P. (1967). Hume's standard of taste: Breaking the circle. \emph{The British Journal of Aesthetics}, 7(1), 57--66.
\bibitem[Levinson(2001)]{Levinson2001} Levinson, J. (2001). Aesthetic properties, evaluative force, and differences of sensibility. In E. Brady and J. Levinson (Eds.), \emph{Aesthetic Concepts: Essays After Sibley}. Oxford University Press.
\bibitem[Levinson(2002)]{Levinson2002} Levinson, J. (2002). Hume's standard of taste: The real problem. \emph{The Journal of Aesthetics and Art Criticism}, 60(3), 227--238.
\bibitem[Meskin et~al.(2013)]{Meskin2013} Meskin, A., Phelan, M., Moore, M. and Kieran, M. (2013). Mere exposure to bad art. \emph{The British Journal of Aesthetics}, 53(2), 139--164.
\bibitem[Mothersill(1984)]{Mothersill1984} Mothersill, M. (1984). \emph{Beauty Restored}. Clarendon Press.
\bibitem[Robson(2012)]{Robson2012} Robson, J. (2012). Aesthetic testimony. \emph{Philosophy Compass}, 7(1), 1--10.
\bibitem[Ross(2008)]{Ross2008} Ross, S. (2008). Humean critics: Real or ideal? \emph{The British Journal of Aesthetics}, 48(1), 20--28.
\bibitem[Shelley(1994)]{Shelley1994} Shelley, J. (1994). Hume's double standard of taste. \emph{The Journal of Aesthetics and Art Criticism}, 52(4), 437--445.
\bibitem[Shelley(2013)]{Shelley2013} Shelley, J. (2013). Hume and the joint verdict of true judges. \emph{The Journal of Aesthetics and Art Criticism}, 71(2), 145--153.
\bibitem[Sibley(1959)]{Sibley1959} Sibley, F. (1959). Aesthetic concepts. \emph{The Philosophical Review}, 68(4), 421--450.
\bibitem[Vessel et~al.(2018)]{Vessel2018} Vessel, E.~A., Maurer, N., Denker, A.~H. and Starr, G.~G. (2018). Stronger shared taste for natural aesthetic domains than for artifacts of human culture. \emph{Cognition}, 179, 121--131.
\bibitem[Walton(1970)]{Walton1970} Walton, K.~L. (1970). Categories of art. \emph{The Philosophical Review}, 79(3), 334--367.
\end{thebibliography}

\end{document}
